\Needspace{20\baselineskip}
\section{Relative moments}
\label{sec:big-relative-first-moment}

For $w\in\mathcal E_\varphi$ and $u\leq w$, the relative first moment $\int_X(w-u)\,\mu_u$ measures how far $u$ lies below $w$, weighted by the measure of $u$. A uniform bound on these moments prevents a sequence from collapsing and keeps its limits in $\mathcal E_\varphi$. The proof compares rooftop masses with the fixed measure of a potential of the same type as $w$ whose mass is close to the upper volume $M$. For a real number $t$, write $t_+=\max(t,0)$.

\begin{proposition}\label{prop:big-relative-first-moment}\index{Relative full-mass class!first-moment criterion}
Fix $w\in\mathcal E_\varphi$. Suppose $u_j\in\mathcal E_\varphi$, $u_j\leq w$, and
\[
 K:=\sup_j\int_X(w-u_j)\,\mu_{u_j}<+\infty.
\]
Then $\sup_Xu_j$ is bounded below. If $u_j\to u\in\PSH(X,\theta)$ in $L^1(X)$, then $u\in\mathcal E_\varphi$.

The same conclusions hold if $u_j\leq w$ is replaced by a common upper bound $u_j\leq\varphi+C$ and the integral is replaced by $\int_X(w-u_j)_+\,\mu_{u_j}$.
\end{proposition}
\begin{proof}
We first treat the case $u_j\leq w$. Since $\overline{\operatorname{vol}}_\theta(w)=M$, we may choose $w'\sim w$ with $\mu_{w'}(X)>M-\ell/2$. Translate $w'$ so that $w'\leq w$, and put $p_j=P_\theta(u_j,w')$. These rooftops belong to $\mathcal E_\varphi$. Contact support and ordered contact imply
\begin{equation}\label{eq:big-energy-roof-bound}
 \int_X(w'-p_j)\,\mu_{p_j}
 \leq\int_X(w-u_j)\,\mu_{u_j}\leq K.
\end{equation}
Indeed, the integrand vanishes on the contact part $\{p_j=w'\}$, and on the remaining contact part $p_j=u_j<w'\leq w$.

\Needspace{8\baselineskip}
We shall use the following estimate for maxima. If $a\in\mathcal E_\varphi$, $a\leq w'$, and $t>0$, then
\begin{equation}\label{eq:big-energy-max-bound}
 \mu_a(\{a>w'-t\})
 \leq \mu_{w'}(\{a>w'-t\})+M-\mu_{w'}(X).
\end{equation}
To see this, put $m=\max(a,w'-t)$. It has the ordinary singularity type of $w'$, so $\mu_m(X)\leq M$. On $\{a>w'-t\}$, locality gives $\mu_m=\mu_a$. On the complementary set, where $m=w'-t$, ordered contact gives $\mu_m\geq\mu_{w'}$. Integrating over the two sets proves \eqref{eq:big-energy-max-bound}.

Suppose that a subsequence had $\sup_Xu_j\to-\infty$; then also $\sup_Xp_j\to-\infty$. Splitting the mass of $p_j$ at the level $w'-s$, $s>0$, and using \eqref{eq:big-energy-roof-bound} and \eqref{eq:big-energy-max-bound}, we obtain
\[
 \ell\leq\mu_{p_j}(X)
 \leq\frac Ks+\mu_{w'}(\{p_j>w'-s\})+M-\mu_{w'}(X).
\]
For fixed $s$, the middle term tends to zero as $j\to\infty$, since $w'$ is finite off a pluripolar set and $\mu_{w'}$ is finite and non-pluripolar. Letting $j\to\infty$ and then $s\to\infty$ gives $\ell\leq M-\mu_{w'}(X)<\ell/2$, a contradiction. This proves the first assertion.

\Needspace{8\baselineskip}
Now suppose $u_j\to u$ in $L^1$. Then $u\leq w\preceq\varphi$. Fix $b>1$ and suppose that $Q_b(u,w)\equiv-\infty$. Since $w-w'$ is bounded, translating competitors by a constant shows that $Q_b(u,w')\equiv-\infty$ as well. Keep the same $p_j$ and set
\[
 h_j=Q_b(p_j,w')\in\mathcal E_\varphi.
\]
We claim that $\sup_Xh_j\to-\infty$. Otherwise, since $h_j\leq p_j\leq w'$, \cref{thm:normalized-compactness} would give a subsequence converging in $L^1$ to some $h\not\equiv-\infty$. Passing to the limit in
\[
 b^{-1}h_j+(1-b^{-1})w'\leq p_j\leq u_j,
\]
first almost everywhere and then everywhere by \cref{lem:qpsh-almost-everywhere-order}, would give a competitor for $Q_b(u,w')$, a contradiction.

On the contact set of $h_j$ one has
$p_j=b^{-1}h_j+(1-b^{-1})w'$. Thus $h_j\leq w'-bs$ implies $p_j\leq w'-s$ there. Contact support and \eqref{eq:big-energy-roof-bound} therefore imply
\[
 \mu_{h_j}(\{h_j\leq w'-bs\})\leq b^nK/s.
\]
Applying \eqref{eq:big-energy-max-bound} with $a=h_j$ and $t=bs$ on the complementary set, we obtain
\[
 \ell\leq b^nK/s+
     \mu_{w'}(\{h_j>w'-bs\})+M-\mu_{w'}(X).
\]
\Needspace{7\baselineskip}
Let first $j\to\infty$ and then $s\to\infty$. The choice of $w'$ again gives $\ell<\ell/2$, a contradiction. Therefore $Q_b(u,w)\not\equiv-\infty$ for every $b>1$. The resulting convex inequalities and \cref{thm:volume-monotonicity,prop:volume-addition} give
\[
 (1-b^{-1})^n M\leq\overline{\operatorname{vol}}_\theta(u)\leq M.
\]
Letting $b\to\infty$ shows that the upper volumes of $u$ and $\varphi$ are equal; hence $u\in\mathcal E_\varphi$.

\Needspace{6\baselineskip}
For the last assertion put $q_j=P_\theta(u_j,w)$. Contact gives
\[
 \int_X(w-q_j)\,\mu_{q_j}
 \leq\int_X(w-u_j)_+\,\mu_{u_j}\leq K.
\]
The case already treated bounds $\sup_Xq_j$ from below, and therefore also $\sup_Xu_j\geq\sup_Xq_j$. If $u_j\to u$ in $L^1$, then, since $q_j\leq w$, a subsequence of the $q_j$ converges in $L^1$ to some $q\not\equiv-\infty$, and the case already treated gives $q\in\mathcal E_\varphi$. Since $q\leq u\leq\varphi+C$, upward stability gives $u\in\mathcal E_\varphi$ as well.
\end{proof}

A common lower bound in $\mathcal E_\varphi$ and a weighted first moment of the difference from the limit tending to zero give convergence in capacity: rooftops of the tails identify the limit.

\Needspace{13\baselineskip}
\begin{lemma}\label{lem:big-weighted-common-bound}
Let $a_j\in\PSH(X,\theta)$ have a common upper bound and converge in $L^1(X)$ to $a\in\PSH(X,\theta)$ with $a\preceq\varphi$. Suppose there is $h\in\mathcal E_\varphi$ such that $h\leq a_j$ for every $j$. If, for some $\lambda>0$,
\[
 \int_X|a_j-a|\mathrm{e}^{\lambda a}\,\mu_{a_j}\longrightarrow0,
\]
then $a\in\mathcal E_\varphi$ and $a_j\to a$ in capacity.
\end{lemma}
\begin{proof}
By \cref{lem:qpsh-almost-everywhere-order}, the lower bound passes to the limit, so $h\leq a\preceq\varphi$, and upward stability gives $a\in\mathcal E_\varphi$. Start with any subsequence and extract a further subsequence along which the integrals in the hypothesis are at most $2^{-j}$ and $a_j\to a$ almost everywhere for smooth volume. For $N\geq j$ let $p_{j,N}=P_\theta(a_j,\ldots,a_N)$; these rooftops are not identically $-\infty$ because $h$ is a competitor. Partitioning the contact set according to the first obstacle touched gives
\[
 \int_X\min(1,|p_{j,N}-a|)\mathrm{e}^{\lambda a}\,\mu_{p_{j,N}}
 \leq\sum_{k=j}^N\int_X|a_k-a|\mathrm{e}^{\lambda a}\,\mu_{a_k}
 \leq2^{1-j}.
\]
The decreasing limits $p_j=P_\theta(\inf_{k\geq j}a_k)$ satisfy $h\leq p_j\leq a$, the second inequality by the almost-everywhere convergence, and hence so does their increasing regularized limit $p$. Upward stability puts the $p_j$ and $p$ in $\mathcal E_\varphi$, although the finite rooftops need not lie there.

Both $p_{j,N}\downarrow p_j$ and $p_j\uparrow p$ hold in capacity by \cref{lem:monotone-psh-capacity}, since $h$ is a common lower bound. Apply \cref{lem:weighted-extended-lsc} to this rooftop estimate, with the bounded weights $\min(1,|p_{j,N}-a|)\mathrm{e}^{\lambda a}$ and then $\min(1,|p_j-a|)\mathrm{e}^{\lambda a}$. These weights converge in capacity: discard the deep sublevel sets of $h$ and $a$, and use continuity on the remaining finite ranges. The result is
\[
 \int_X\min(1,|p-a|)\mathrm{e}^{\lambda a}\,\mu_p=0.
\]
Since the exponential is positive off a pluripolar set, which $\mu_p$ does not charge, $\mu_p(\{p<a\})=0$, and zero domination gives $p=a$. The upper envelopes $(\sup_{k\geq j}a_k)^*$ decrease to $a$. Squeezing the extracted $a_j$ between these envelopes and $p_j$ proves convergence in capacity. Since the initial subsequence was arbitrary, the conclusion holds for the original sequence.
\end{proof}

Without a given lower bound, the first-moment criterion produces one from the rooftops of the tails.

\Needspace{20\baselineskip}
\begin{theorem}\label{thm:big-relative-self-measure-convergence}
Let $u_j,u\in\mathcal E_\varphi$, with $u_j\to u$ in $L^1(X)$. If
\[
 \int_X|u_j-u|\,\mu_{u_j}\longrightarrow0,
\]
then $u_j\to u$ in capacity. Moreover, after passing to a subsequence, $P_\theta(\inf_{k\geq j}u_k)$ belongs to $\mathcal E_\varphi$ for every $j$.
\end{theorem}
\begin{proof}
Start with any subsequence and extract a further subsequence along which the integrals in the hypothesis are at most $2^{-j}$ and $u_j\to u$ almost everywhere for smooth volume. The envelope identity and Hartogs' lemma give a uniform bound $u_j,u\leq\varphi+C$. For $N\geq j$ let $p_{j,N}=P_\theta(u_j,\ldots,u_N)$. Partitioning the contact set as in the proof of the preceding lemma gives
\begin{equation}\label{eq:big-self-measure-roofs}
 \int_X|p_{j,N}-u|\,\mu_{p_{j,N}}\leq2^{1-j}.
\end{equation}
By \eqref{eq:big-self-measure-roofs}, \cref{prop:big-relative-first-moment}, applied with $w=u$ in its one-sided form, bounds $\sup_Xp_{j,N}$ from below, independently of $N$. Consequently the decreasing sequence has a limit
\[
 p_j=P_\theta\left(\inf_{k\geq j}u_k\right)
\]
that is not identically $-\infty$, and the same proposition puts $p_j$ in $\mathcal E_\varphi$. This proves the second assertion. Since $p_1\leq u_k$ for every $k$ and $\mathrm{e}^{u}\leq\mathrm{e}^{\sup_Xu}$, the preceding lemma, applied to the extracted sequence with $h=p_1$ and $\lambda=1$, shows that it converges to $u$ in capacity. Since the initial subsequence was arbitrary, the conclusion holds for the original sequence.
\end{proof}
